\section{Results and Discussion}

\newrev{MBA-RL was evaluated using SPMe models identified from three Samsung 40T 21700 cells, each with a nominal capacity of 4000 mAh. Voltage and temperature measurements were collected using the NEWARE platform in Fig.~\ref{fig:experiment_Schematic}. The identified models were used for training and simulation tests with independently regulated cell charging currents.}

\begin{figure}[t]
    \centering
    \includegraphics[width=\columnwidth]{experimental_platform.png}
    \caption{Experimental platform for cell voltage and temperature measurements.}
    \label{fig:experiment_Schematic}
\end{figure}

\subsection{Cell Model Identification}

\newrev{Cells 2 and 3 underwent 140 and 260 charge-discharge cycles, respectively, to obtain cells with different cycling histories. Their voltage and temperature responses, together with those of Cell 1, are shown in Fig.~\ref{fig:three_aging}. Cell-specific SPMe parameters were identified using particle swarm optimization (PSO) \cite{28}. Figs.~\ref{battery1}--\ref{battery3} compare the model predictions with the measurements. The voltage and temperature RMSEs in Table~\ref{tab:model_rmse} quantify the fit under these identification conditions.}

\begin{figure}[t]
    \centering
    \begin{subfigure}[t]{0.48\columnwidth}
        \includegraphics[width=\linewidth]{aging_voltage.png}
        \caption{}
    \end{subfigure}\hfill
    \begin{subfigure}[t]{0.48\columnwidth}
        \includegraphics[width=\linewidth]{aging_temperature.png}
        \caption{}
    \end{subfigure}
    \caption{Measured responses of the three cells: (a) voltage; (b) temperature.}
    \label{fig:three_aging}
\end{figure}

\begin{figure}[t]
    \centering
    \begin{subfigure}[t]{0.48\columnwidth}
        \includegraphics[width=\linewidth]{cell1_voltage_fit.png}
        \caption{}
    \end{subfigure}\hfill
    \begin{subfigure}[t]{0.48\columnwidth}
        \includegraphics[width=\linewidth]{cell1_temperature_fit.png}
        \caption{}
    \end{subfigure}
    \caption{Measured and modeled responses of Cell 1: (a) voltage; (b) temperature.}
    \label{battery1}
\end{figure}

\begin{figure}[t]
    \centering
    \begin{subfigure}[t]{0.48\columnwidth}
        \includegraphics[width=\linewidth]{cell2_voltage_fit.png}
        \caption{}
    \end{subfigure}\hfill
    \begin{subfigure}[t]{0.48\columnwidth}
        \includegraphics[width=\linewidth]{cell2_temperature_fit.png}
        \caption{}
    \end{subfigure}
    \caption{Measured and modeled responses of Cell 2: (a) voltage; (b) temperature.}
    \label{battery2}
\end{figure}

\begin{figure}[t]
    \centering
    \begin{subfigure}[t]{0.48\columnwidth}
        \includegraphics[width=\linewidth]{cell3_voltage_fit.png}
        \caption{}
    \end{subfigure}\hfill
    \begin{subfigure}[t]{0.48\columnwidth}
        \includegraphics[width=\linewidth]{cell3_temperature_fit.png}
        \caption{}
    \end{subfigure}
    \caption{Measured and modeled responses of Cell 3: (a) voltage; (b) temperature.}
    \label{battery3}
\end{figure}

\begin{table}[t]
    \centering
    \caption{Voltage and temperature RMSEs of the identified cell models}
    \label{tab:model_rmse}
    \begin{tabular}{lccc}
        \toprule
        Metric & Cell 1 & Cell 2 & Cell 3 \\
        \midrule
        Voltage RMSE (V) & 0.026 & 0.035 & 0.029 \\
        Temperature RMSE (K) & 0.068 & 0.091 & 0.093 \\
        \bottomrule
    \end{tabular}
\end{table}

\subsection{Evaluation Protocol}

\newrev{The cell models were implemented in PyBaMM \cite{29}. Policies were trained through interaction with these models and evaluated with fixed network parameters. The control and evaluation settings are summarized in Table~\ref{tab:experimental_settings}. Charging ends when every cell reaches $SOC_{\mathrm{ref}}=0.90$ and $\sigma(t)\leq\beta=0.02$. A subsequent 270 s verification period at zero current checks that the SOC target, balancing criterion, and safety limits remain satisfied. A test is successful only if the charging target and this verification are completed without a constraint violation at the sampled times.}

\begin{table}[t]
    \centering
    \color{red}
    \caption{Control and evaluation settings}
    \label{tab:experimental_settings}
    \footnotesize
    \renewcommand{\arraystretch}{1.08}
    \begin{tabular}{@{}p{0.35\columnwidth}p{0.59\columnwidth}@{}}
        \toprule
        Setting & Value \\
        \midrule
        Cell model & SPMe with lumped thermal dynamics \\
        Discrete current set & $0$--$7.5$ A in steps of 0.5 A \\
        Continuous current range & $0$--$7.5$ A \\
        Decision interval & 90 s \\
        Charging target & $SOC_i\geq0.90$ for all cells; $\sigma\leq0.02$ \\
        Maximum charging time & 5400 s \\
        Verification period & 270 s at zero current \\
        Operating limits & $V\leq4.2$ V, $T\leq309$ K, $SOC\leq0.95$ \\
        Voltage / temperature margins & 0.03 V / 1 K \\
        Ambient temperature & 298.15 K (25~$^\circ$C) \\
        \bottomrule
    \end{tabular}
\end{table}

\newrev{Performance is assessed by task completion, constraint satisfaction, a time score that penalizes unsuccessful tests, and accumulated SOC imbalance. For a successful test, the time score is the charging time $t_b-t_0$, excluding the verification period. An unsuccessful test receives a score of $5400+270=5670$ s. Accumulated SOC imbalance is defined as}
\begin{equation}
    \label{eq:imbalance_metric}
    J_{\sigma}=\int_{t_0}^{t_{\mathrm{end}}}\sigma(t)\,\mathrm{d}t,
\end{equation}
\newrev{where $t_{\mathrm{end}}$ is the last recorded time in the charging phase. The integral excludes the zero-current verification period and uses the recorded charging trajectory even if the test is unsuccessful; it is not extended to the failure time score. SOC is expressed as a fraction, so $J_{\sigma}$ is reported in SOC fraction\,$\cdot$\,s.}

\subsection{Charging Performance}

\subsubsection{Illustrative Cell Trajectories}

\newrev{Fig.~\ref{fig:overall_figure} shows the cell SOC trajectories for initial conditions Set A, $[0.10,0.20,0.30]^{\mathrm{T}}$, and Set B, $[0.30,0.50,0.70]^{\mathrm{T}}$. MBA-RL selects each cell's charging current from its local observations, with coordination learned during centralized training. The DQL controller in these illustrative comparisons also controls a cell-to-cell energy-transfer topology, whereas the CC-CV baseline activates that topology during the constant-voltage phase.}

\begin{figure*}[t]
    \centering
    \begin{subfigure}[t]{0.32\textwidth}
        \includegraphics[width=\linewidth]{mba_rl_set_a.png}
        \caption{}
    \end{subfigure}\hfill
    \begin{subfigure}[t]{0.32\textwidth}
        \includegraphics[width=\linewidth]{dql_set_a.png}
        \caption{}
    \end{subfigure}\hfill
    \begin{subfigure}[t]{0.32\textwidth}
        \includegraphics[width=\linewidth]{cc_cv_set_a.png}
        \caption{}
    \end{subfigure}

    \medskip
    \begin{subfigure}[t]{0.32\textwidth}
        \includegraphics[width=\linewidth]{mba_rl_set_b.png}
        \caption{}
    \end{subfigure}\hfill
    \begin{subfigure}[t]{0.32\textwidth}
        \includegraphics[width=\linewidth]{dql_set_b.png}
        \caption{}
    \end{subfigure}\hfill
    \begin{subfigure}[t]{0.32\textwidth}
        \includegraphics[width=\linewidth]{cc_cv_set_b.png}
        \caption{}
    \end{subfigure}
    \caption{\newrev{Cell SOC trajectories for MBA-RL, DQL, and CC-CV (left to right): (a)--(c) Set A; (d)--(f) Set B. Green dashed lines mark the first time $\sigma\leq\beta$; yellow dashed lines mark the end of constant-current charging.}}
    \label{fig:overall_figure}
\end{figure*}

\newrev{All three methods first meet the balancing threshold before every cell reaches the target SOC. The MBA-RL trajectories continue after this point to illustrate subsequent charging, whereas the displayed DQL and CC-CV trajectories end at balance. Table~\ref{tab:soc_vectors} lists the SOC vectors at the ends of these displayed trajectories. Their different endpoints distinguish initial balance from completion of the charging task; quantitative comparisons using a common task criterion are presented next.}

\begin{table}[t]
    \centering
    \caption{\newrev{Cell SOCs at the ends of the illustrative trajectories}}
    \label{tab:soc_vectors}
    \footnotesize
    \setlength{\tabcolsep}{4pt}
    \begin{tabular}{@{}lcc@{}}
        \toprule
        Method & Initial SOC (\%) & Final plotted SOC (\%) \\
        \midrule
        \multirow{2}{*}{MBA-RL} & $[10,20,30]^{\mathrm{T}}$ & $[91.66,94.18,92.15]^{\mathrm{T}}$ \\
         & $[30,50,70]^{\mathrm{T}}$ & $[91.23,93.62,91.41]^{\mathrm{T}}$ \\
        \addlinespace
        \multirow{2}{*}{DQL} & $[10,20,30]^{\mathrm{T}}$ & $[24.72,26.81,28.05]^{\mathrm{T}}$ \\
         & $[30,50,70]^{\mathrm{T}}$ & $[66.50,69.18,70.00]^{\mathrm{T}}$ \\
        \addlinespace
        \multirow{2}{*}{CC-CV} & $[10,20,30]^{\mathrm{T}}$ & $[52.41,55.42,56.97]^{\mathrm{T}}$ \\
         & $[30,50,70]^{\mathrm{T}}$ & $[76.74,79.21,81.38]^{\mathrm{T}}$ \\
        \bottomrule
    \end{tabular}
\end{table}

\subsubsection{Comparison under Unseen Initial SOC Conditions}

\newrev{MBA-RL was compared with parameter-sharing independent Q-learning (shared IQL), centralized DQN \cite{mnih2015dqn}, discrete and continuous MAPPO \cite{yu2022mappo}, and SAC \cite{haarnoja2018sac}. The centralized DQN baseline directly selects cell charging currents and is distinct from the topology-based DQL controller in Fig.~\ref{fig:overall_figure}. The same operating limits and model-based safety screening were applied during training and evaluation. Each method was trained independently five times and evaluated under the same 12 initial SOC conditions not used during training, giving 60 tests per method.}

\begin{table}[t]
    \centering
    \color{red}
    \caption{Performance comparison under unseen initial SOC conditions}
    \label{tab:modern_baselines}
    \footnotesize
    \renewcommand{\arraystretch}{1.1}
    \begin{tabular*}{\columnwidth}{@{\extracolsep{\fill}}lcc@{}}
        \toprule
        Method & Time score (s) & $J_{\sigma}$ \\
         & & (SOC fraction\,$\cdot$\,s) \\
        \midrule
        \multicolumn{3}{@{}l}{\textit{Discrete-action methods}} \\
        MBA-RL & $2701.50\pm919.80$ & $205.42\pm57.03$ \\
        Shared IQL & $2641.50\pm439.42$ & $235.17\pm69.76$ \\
        Centralized DQN & $2673.00\pm73.98$ & $230.35\pm24.20$ \\
        Discrete MAPPO & $2340.00\pm347.96$ & $200.83\pm60.75$ \\
        \addlinespace
        \multicolumn{3}{@{}l}{\textit{Continuous-action methods}} \\
        Continuous MAPPO & $2698.50\pm154.65$ & $240.36\pm27.58$ \\
        SAC & $2173.50\pm157.09$ & $156.04\pm30.74$ \\
        \bottomrule
    \end{tabular*}
    \par\smallskip
    \parbox{\columnwidth}{\scriptsize Values are means $\pm$ standard deviations across five independently trained policies; each policy value is averaged over 12 test conditions.}
\end{table}

\newrev{All methods completed 60/60 tests except continuous MAPPO, which completed 59/60. Under the same discrete current set, MBA-RL reduced accumulated SOC imbalance by 12.65\% relative to shared IQL and by 10.82\% relative to centralized DQN, with a slightly higher mean time score. Discrete MAPPO and SAC achieved lower mean time scores and accumulated SOC imbalance than MBA-RL. Thus, the improvement over the two value-based baselines concerns SOC balancing rather than shorter charging time.}

\newrev{Fig.~\ref{fig:discrete_trajectories} illustrates the responses for initial SOCs $[0.4203,0.3962,0.5280]^{\mathrm{T}}$ at the ambient temperature in Table~\ref{tab:experimental_settings}. MBA-RL reaches a mean SOC of 0.90 earlier than shared IQL and has a lower final SOC standard deviation than centralized DQN in this case. Task completion is assessed using the individual cell SOCs and the verification period defined above.}

\begin{figure}[t]
    \centering
    \includegraphics[width=\columnwidth]{comparative_charging_discrete.pdf}
    \caption{\newrev{Charging trajectories of MBA-RL, shared IQL, and centralized DQN: (a) mean SOC; (b) SOC standard deviation. The dashed line in (b) marks $\beta=0.02$. The zero-current verification period is not shown.}}
    \label{fig:discrete_trajectories}
\end{figure}

\subsection{Effect of the Safety Mechanism}

\newrev{The safety mechanism was evaluated in three settings: no safety mechanism, screening only during execution of a fixed policy, and screening during both training and execution. Each setting was tested under 26 conditions for two independently trained policies. Table~\ref{tab:safety_results} reports constraint satisfaction at the sampled times and task completion.}

\begin{table}[t]
    \centering
    \color{red}
    \caption{Effect of safety screening on constraint satisfaction and task completion}
    \label{tab:safety_results}
    \footnotesize
    \begin{tabular*}{\columnwidth}{@{\extracolsep{\fill}}lcc@{}}
        \toprule
        Screening setting & Constraints satisfied & Tasks completed \\
        \midrule
        None & 3/52 & 1/52 \\
        Execution only & 52/52 & 35/52 \\
        Training and execution & 52/52 & 52/52 \\
        \bottomrule
    \end{tabular*}
\end{table}

\newrev{Execution screening increased constraint satisfaction from 3/52 to 52/52, but 17 tests still failed to complete charging. Including the mechanism during training increased task completion to 52/52 and reduced mean accumulated SOC imbalance from 469.97 to 197.34 SOC fraction\,$\cdot$\,s compared with execution-only screening. Training with the mechanism therefore improved charging performance under the screened action constraints.}

\subsection{Component Ablation}

\newrev{The contributions of parameter sharing, the mixing network, and the balancing reward were examined using the same evaluation protocol as the main comparison.}

\begin{figure}[t]
    \centering
    \includegraphics[width=\columnwidth]{component_ablation.pdf}
    \caption{\newrev{Component ablation results. Circles denote each trained policy's mean over 12 test conditions; diamonds denote the means across policies. The dashed line marks the MBA-RL mean; right-hand labels report completed tests out of 60.}}
    \label{fig:component_ablation}
\end{figure}

\newrev{Removing parameter sharing or the balancing reward reduced task completion from 60/60 to 48/60 and increased mean accumulated SOC imbalance by 80.47\% and 59.07\%, respectively. In each case, one of the five trained policies failed all 12 test conditions. Removing the mixing network retained 60/60 completion but increased accumulated SOC imbalance by 14.48\%. These results show that parameter sharing and the balancing reward improve the consistency of performance across independently trained policies in the three-cell task.}

\subsection{Effects of Operating Conditions}

\newrev{The trained policy was evaluated under parameter variations, sensor errors, observation delay, ambient-temperature changes, and reduced current limits. Each condition comprised 100 tests, including 50 with initially balanced cells and 50 with random initial SOCs. All tests were completed under the examined $\pm10\%$ parameter variations, fixed sensor bias, 15~$^\circ$C ambient temperature, and a 30\% reduction in the current limit. Completion decreased with observation noise, a one-decision observation delay, and an ambient temperature of 34~$^\circ$C. At 35~$^\circ$C, which was included to test feasibility under the prescribed safety margins, no test completed charging. The results therefore indicate greater sensitivity to observation errors and high ambient temperature than to the tested parameter and current-limit changes.}

\subsection{Effect of Pack Size}

\newrev{The architecture was trained separately for 3, 6, and 12 independently controlled cells. These scale tests used separately generated initial SOC conditions, so their three-cell results differ from those in Table~\ref{tab:modern_baselines}. Task completion decreased from 59/60 for three cells to 29/36 for twelve cells. The twelve-cell result is preliminary and comprises the three training runs completed at the time of analysis; five runs were completed for the three-cell case.}

\newrev{Policy inference required approximately 2--3 ms, while safety screening increased from about 1.7 s for three cells to 7.5 s for twelve cells. Both remained below the 90 s decision interval in these tests. Computational time was therefore compatible with the control interval, although charging performance declined as the number of cells increased. The larger simulated packs repeat the three identified cell models and do not include series-circuit coupling, a shared power limit, or heat transfer between cells.}
